\begin{figure}[ht]
  \centering
  \begin{tikzpicture}[
      box/.style={draw, rounded corners=2pt, align=center, font=\footnotesize, inner sep=4pt, minimum height=8mm, text width=4.6cm},
      tratt/.style={box, dashed},
      freccia/.style={-{Stealth[length=5pt]}, line width=.7pt}]
    \node[box, text width=3cm] (q) at (0,0) {query};
    \node[box] (lex) at (-2.7,-1.5) {recupero lessicale\\ disgiuntivo o congiuntivo};
    \node[tratt] (vec) at (2.7,-1.5) {i 100 documenti più vicini\\ per \texttt{bge-m3}};
    \node[tratt, text width=6cm] (uni) at (0,-3.1) {unione dei candidati};
    \node[tratt, text width=6cm] (fus) at (0,-4.3) {fusione dei punteggi, con il peso del vettore};
    \node[tratt, text width=6cm] (rr) at (0,-5.5) {cross-encoder sui primi $k$ candidati};
    \node[box, text width=3cm] (ris) at (0,-6.7) {risultati};
    \draw[freccia] (q) -- (lex); \draw[freccia] (q) -- (vec);
    \draw[freccia] (lex) -- (uni); \draw[freccia] (vec) -- (uni);
    \draw[freccia] (uni) -- (fus); \draw[freccia] (fus) -- (rr); \draw[freccia] (rr) -- (ris);
  \end{tikzpicture}
  \caption[Le configurazioni a confronto, passo per passo]{Dove differiscono le configurazioni del capitolo. Il tratteggio segna
  i passi che alcune non fanno. Il recupero lessicale è disgiuntivo in LA, A e
  P2, P3, congiuntivo in LT, B e P4. Vettori e unione ci sono in A, B, P3 e P4,
  con peso 160 in A e P3 e peso 10 in B e P4. Il cross-encoder, fuori da
  Koskidex, riordina i primi $k$ candidati: 100 in RR-LA e RR-A, 20 o 30 nel
  riordino corto.}
  \label{fig:flusso}
\end{figure}
